% =====================================================================
% DRAFT — SQ2: does the projection view (plane A vs B) change the AIF and the perfusion parameters?
% Methods -> sections/02_methods.tex, Results -> sections/03_results.tex,
% Discussion -> sections/04_discussion.tex. References: methods_aif_refs.bib.
%
% Show in the main text:
%   Table~\ref{tab:sq2_main}   results/combined_v1/sq2_plane/table_sq2.tex (ica_top, 13 features)
%   Fig.~\ref{fig:sq2_curves}  sq2_plane/shape_curves.png -> figures/sq2_shape_curves.png
%   Fig.~\ref{fig:sq2_mech}    sq2_plane/mechanism.png    -> figures/sq2_mechanism.png
%   Fig.~\ref{fig:sq2_forest}  sq2_plane/forest.png       -> figures/sq2_forest.png
% Appendix D: ba_ica_top.png, ba_ica_full_model.png, ba_ica_full_grown.png, qq_*.png, sq2_report.md tables.
%
% Numbers: results/combined_v1/sq2_plane/sq2_report.md (QC with tail threshold 0.30, current pipeline).
% =====================================================================

% ---------------------------------------------------------------------
% METHODS
% ---------------------------------------------------------------------
\subsection{Projection view (SQ2)}
\label{sec:sq2_methods}

\textbf{Rationale.} A DSA image is a projection: each pixel holds the iodine integrated along one X-ray path.
The same bolus therefore produces a different AIF, and different tissue curves, in the frontal and the lateral
view, depending on how long the path through the vessel is, what else lies on that path, and how the image was
exposed. Deconvolution-based perfusion parameters are expected to be less sensitive to such differences than raw
time--intensity features, because a factor that scales the AIF and the tissue curves equally cancels in the
deconvolution; in an animal study they varied two to five times less with the injection
protocol~\cite{su2023animal}. SQ2 tests whether this holds between the two views of one injection.

\textbf{Pairing.} The biplane system records both planes of an injection simultaneously with identical frame
times, so an A/B pair shares the bolus, the patient state and the timing; any difference is introduced by the
projection. Pairs were formed per injection and AIF variant when both curves were trusted
(Section~\ref{sec:aif_qc}). The primary analysis uses \texttt{ica\_top}, the AIF recommended by SQ1; the
full-ICA variants are reported for completeness. No patient contributed more than one \texttt{ica\_top} pair,
so the pairs are independent.

\textbf{Paired comparison and agreement.} Differences B $-$ A were computed as in SQ1: plain differences for
times and dimensionless shape features, log ratios B/A for amplitudes and perfusion parameters, whose spread
grows with their size~\cite{bland1999}. Normality was assessed with Shapiro--Wilk and Q--Q plots; the paired
$t$-test was used for normal differences and the Wilcoxon signed-rank test otherwise, with Holm correction over
the 15 features of one variant. Agreement was quantified with Bland--Altman 95\,\% limits of
agreement~\cite{bland1999}, parametric (bias $\pm$ 1.96\,SD) for normal and non-parametric (2.5th--97.5th
percentile) for non-normal differences.

\textbf{Why the AIF differs between views.} Four explanations for a higher AIF in one view were made testable:
(i) \emph{exposure}: the automatic exposure control sets a lower tube voltage for the thinner lateral path through
the head, and iodine attenuates more at lower tube voltage~\cite{cho2012}; this raises all contrast in the image,
so the AIF ratio B/A should equal the whole-image contrast ratio B/A, i.e.\ AIF normalised to whole-image contrast
should be equal in both views; (ii) \emph{path length}: a foreshortened siphon projects onto fewer pixels with a
longer path per pixel, so the AIF ratio would relate to the mask-size ratio; (iii) \emph{saturation}: clipping of
the densest pixels can only lower the measured peak, so pairs with saturated mask pixels should show a lower
ratio; (iv) \emph{overlap}: vessels projected over the mask add iodine, measured as the vessel fraction in a
15-pixel ring around the mask. Each was related to the log AIF peak ratio with Spearman correlation, and pairs
with and without saturated pixels were compared with the Mann--Whitney test.

\textbf{Does normalisation remove the view difference?} If the views differ mainly by a scale factor, removing
the amplitude should restore agreement. This was tested on features (bias and limits of agreement of the
shape-only features against those of the amplitude features) and on whole curves: the A and B curves were
aligned at their arrival times and the root-mean-square difference between them (1\,s before to 8\,s after
arrival, in units of the mean peak) was computed for the raw, the peak-normalised and the AUC-normalised curves,
and compared with the paired Wilcoxon test.

\textbf{Sensitivity analyses.} The main comparisons were repeated (i) for ICA injections only, because a CCA
injection fills ECA branches that overlap the ICA differently in each view, and (ii) without the pairs that had
saturated mask pixels.

% ---------------------------------------------------------------------
% RESULTS
% ---------------------------------------------------------------------
\subsection{Projection view (SQ2)}
\label{sec:sq2_results}

\textbf{Pairs and acquisition.} Both planes were trusted in 43 injections for \texttt{ica\_top} (42
\texttt{ica\_full\_model}, 57 \texttt{ica\_full\_grown}), from 43 different patients. Frame times, field of view,
magnification (1.46 in both) and detector were identical between the views; the exposure was not: plane B was
acquired at 69.2 vs.\ 76.8\,kVp with about a third of the dose--area product. The AIF curves of the two views were
highly correlated in time (median $r$ = 0.96). Shapiro--Wilk rejected normality for 11 of 15 \texttt{ica\_top}
features; the $t$-test and Wilcoxon test disagreed on significance in 4 of 45 tests (\texttt{ica\_top} AUC/peak
and CBF, and CBV of both full-ICA variants), so those four results are reported as direction only.

\textbf{The AIF is higher in the lateral view.} The \texttt{ica\_top} peak was 36\,\% higher and the first-pass AUC
42\,\% higher in plane B (Table~\ref{tab:sq2_main}, Fig.~\ref{fig:sq2_curves}; $p_\mathrm{Holm} < 10^{-8}$), with
limits of agreement of $-1$\,\% to $+87$\,\% for the peak: in almost every injection the lateral AIF was the higher
one. The full-ICA variants showed the same direction (model $+38$\,\%, grown $+12$\,\%; Fig.~\ref{fig:sq2_forest}).
Timing hardly differed: arrival $+0.08$\,s (n.s.), time to peak and Tmax with identical medians, all within one
frame interval.

\textbf{Exposure explains the higher peak; path length and saturation do not.} The whole-image contrast was 32\,\%
higher in plane B, the same factor as the AIF; normalised to it, the \texttt{ica\_top} AIF was equal in both views
(B/A 0.98 in pairs without and 1.03 in pairs with saturated pixels; 1.00 over all 68 A/B pairs, $p_\mathrm{Holm} =
1.0$, Section~\ref{sec:plane_methods}). The AIF ratio increased with the tube-voltage difference ($\rho = -0.46$,
$p = 0.002$: the larger the kVp drop in B, the higher the B peak) and with the vessel overlap around the mask
($\rho = 0.46$, $p = 0.002$), but not with the mask-size ratio, the proxy for path length ($\rho = -0.02$, $p =
0.89$), nor with saturation ($\rho = 0.08$, $p = 0.63$; Fig.~\ref{fig:sq2_mech}). Saturated mask pixels occurred
only in plane B (7 of 43 pairs), and these pairs had the same peak ratio as the others (1.36 vs.\ 1.32, $p =
0.74$).

\textbf{Normalisation removes most of the AIF difference.} Of the shape-only features, AUC/peak, rise time,
skewness and tail fraction did not differ between views; the FWHM was 0.13\,s wider in B ($p_\mathrm{Holm} <
10^{-4}$) and the temporal SD 0.07\,s smaller ($p_\mathrm{Holm} = 0.04$), both about half a frame interval.
At curve level, the difference between the A and B curve fell from a median of 17.7\,\% of the peak for the raw
curves to 6.6\,\% after peak normalisation and 6.4\,\% after AUC normalisation (both $p < 10^{-10}$), a reduction
of almost two thirds (Fig.~\ref{fig:sq2_curves}).

\textbf{The perfusion parameters differ even more, in the opposite direction.} Despite the higher AIF, CBV was
57\,\% lower and MTT 40\,\% lower in plane B ($p_\mathrm{Holm} = 0.010$ and $< 0.001$); CBF was 25\,\% lower but not
significantly after correction (Table~\ref{tab:sq2_main}). All three variants showed lower MTT in B ($-33$ to
$-40$\,\%, significant in each) and lower CBV ($-36$ to $-57$\,\%; for the full-ICA variants not significant after
Holm correction with the Wilcoxon test). The limits of agreement of the perfusion parameters
were very wide (CBV $-89$\,\% to more than $+1000$\,\%, driven by a few runs with extreme map values). Because CBV
is the ratio of the tissue to the AIF area, the higher AIF area and the lower CBV together imply that the tissue
signal itself was lower in the lateral view: per pair, CBV $\times$ AIF first-pass AUC gave a tissue-area ratio
B/A of 0.59 (IQR 0.49--0.89, $p = 0.003$), while vessels and the whole image were 32--42\,\% brighter. The
results were unchanged for ICA injections only ($n$ = 26: peak $+31$\,\%, CBV $-57$\,\%, MTT $-40$\,\%) and without
saturated pairs ($n$ = 36: peak $+35$\,\%, CBV $-53$\,\%, MTT $-35$\,\%).

\begin{table}[t]
\centering
\caption{Plane B (lateral) versus plane A (frontal) of the same injection, \texttt{ica\_top} AIF, pairs with both curves trusted ($n$ = 43). Effect: mean (paired $t$) or median ($^\dagger$Wilcoxon) of B $-$ A; amplitudes and perfusion parameters as ratio B/A $-$ 1. LoA: Bland--Altman 95\,\% limits of agreement (bias $\pm$ 1.96\,SD; $^\dagger$2.5th--97.5th percentile). $p$: Holm-corrected over 15 features. Shape-only features in italics.}
\label{tab:sq2_main}
\footnotesize
\begin{tabular}{lrcr}
\hline
Feature & B vs A & 95\,\% LoA & $p$ \\
\hline
Arrival time (s) & +0.080$^\dagger$ & -0.22 to +0.23 & 0.087 \\
Time to peak (s) & +0.000$^\dagger$ & -0.25 to +0.50 & 0.022 \\
Peak height & +35.7\% & -1\% to +87\% & 2.4e-14 \\
First-pass AUC & +42.3\% & -18\% to +147\% & 3.5e-09 \\
\emph{FWHM (s)} & +0.125$^\dagger$ & -0.10 to +0.51 & 1.6e-05 \\
\emph{Rise time (s)} & +0.063 & -0.38 to +0.50 & 0.301 \\
\emph{AUC / peak (s)} & +0.103 & -0.51 to +0.71 & 0.180 \\
\emph{Temporal SD (s)} & -0.065$^\dagger$ & -0.45 to +0.25 & 0.037 \\
\emph{Tail fraction} & -0.001$^\dagger$ & -0.14 to +0.06 & 0.770 \\
CBV (median) & -57.5\%$^\dagger$ & -89\% to +1662\% & 0.010 \\
CBF (median) & -24.6\%$^\dagger$ & -67\% to +705\% & 0.144 \\
MTT (median) & -40.5\%$^\dagger$ & -85\% to +158\% & 6.7e-04 \\
Tmax (median) (s) & +0.000$^\dagger$ & -1.48 to +0.25 & 0.038 \\
\hline
\end{tabular}
\end{table}
  % results/combined_v1/sq2_plane/table_sq2.tex

\begin{figure}[t]
\centering
\includegraphics[width=\columnwidth]{figures/sq2_shape_curves.png}
\caption{\texttt{ica\_top} AIF of plane A (frontal) and B (lateral) of the same injections ($n$ = 43), aligned at
arrival; median and IQR. Left: baseline-corrected, the lateral AIF is higher. Right: peak-normalised, the curves
nearly coincide.}
\label{fig:sq2_curves}
\end{figure}

\begin{figure*}[t]
\centering
\includegraphics[width=\textwidth]{figures/sq2_mechanism.png}
\caption{AIF peak ratio B/A per injection against the candidate explanations: whole-image contrast ratio
(exposure; dashed: equal ratios), mask-size ratio (path length / foreshortening), vessel fraction around the mask
(overlap) and tube-voltage difference. Red: saturated mask pixels in one of the planes.}
\label{fig:sq2_mech}
\end{figure*}

\begin{figure}[t]
\centering
\includegraphics[width=\columnwidth]{figures/sq2_forest.png}
\caption{Plane B minus plane A for every feature and AIF variant: median difference divided by the SD of the
differences (IQR); ringed: $p_\mathrm{Holm} < 0.05$. Amplitudes are higher and CBV and MTT lower in the lateral view
for all three variants.}
\label{fig:sq2_forest}
\end{figure}

% ---------------------------------------------------------------------
% DISCUSSION
% ---------------------------------------------------------------------
\subsection{Does the projection view matter?}
\label{sec:sq2_discussion}

\textbf{The AIF: a scale difference.} The lateral AIF is about a third higher than the frontal AIF of the same
bolus, but this is a property of the image, not of the artery: the whole image is brighter by the same factor,
and after normalising to it the AIF is the same in both views. The most plausible cause is the exposure. The
automatic exposure control chose about 8\,kVp less for the lateral view, and iodine contrast increases at lower
tube voltage~\cite{cho2012}; the AIF ratio indeed grew with the kVp drop. Two other candidate causes can be
excluded as the main driver: path length (no relation with the projected length of the segment) and saturation
(it occurred only in plane B, can only lower a peak, and did not change the ratio). Vessel overlap, which is much
larger around the supraclinoid segment in the lateral view (ring vessel fraction 0.63 vs.\ 0.35), adds to the
pair-to-pair variation (it correlates with the ratio) but does not explain the systematic shift. Because the
difference is mainly a scale factor, removing the amplitude restores most of the agreement: shape-only features
agree between views apart from a FWHM difference of half a frame, and the curve mismatch drops by two thirds.

\textbf{The perfusion maps: not a scale difference.} If the views differed only by a common scale factor, it
would cancel in the deconvolution and CBV, CBF and MTT would agree~\cite{su2023animal}. They do not: CBV and MTT are
40--57\,\% lower in the lateral view, even though the lateral AIF is higher. The implied tissue signal is about
40\,\% lower in the lateral view while vessels are brighter, so vessels and parenchyma respond differently to the
change of view. The AIF swap of SQ4 (Section~\ref{sec:sq4_results}) separates the two: the CBF difference between the
views is caused entirely by the AIF, whereas for CBV and MTT a tissue part remains (about a quarter lower in the lateral
view with the same AIF). The cause of that tissue part cannot be isolated here; candidate explanations for future
work are: (a) the projected
depth of perfused parenchyma per pixel differs between views (front-to-back through the injected hemisphere
versus across it, with the non-injected hemisphere on the same path); (b) the display processing of the vendor
images (pixel intensity relationship DISP) need not be linear for the faint parenchymal signal; (c) the lower dose
of the lateral view raises the noise of the tissue curves, which the truncated-SVD deconvolution and the clipping
of negative values at zero do not treat symmetrically; (d) the perfused-pixel set over which the median is taken
differs per view. Whatever the cause, perfusion values from the two views of one injection are not
interchangeable.

\textbf{Bias and underlying factors to keep in mind.}
\begin{itemize}
\item \emph{Exposure is not controlled}: tube voltage and dose are set per view by the automatic exposure
control, so view and exposure are confounded; the effect of the view cannot be separated from the effect of kVp.
\item \emph{Selection}: pairs required both curves to be trusted. The lateral view had more vessel overlap but also
more trusted curves (54 vs.\ 48 of 68 for \texttt{ica\_top}), so the excluded pairs are not a random subset.
\item \emph{Saturation} occurs only in the lateral view (7 pairs); it did not change the results but could matter
with a stronger injection.
\item \emph{Overlap and injection site}: overlap is view-dependent and larger with CCA injections (ECA branches);
restricting to ICA injections did not change the results.
\item \emph{Segmentation per view}: the \texttt{ica\_top} mask is larger in the lateral view (1586 vs.\ 1129\,px), so
the two AIFs average different stretches of the supraclinoid ICA.
\item \emph{No ground truth}: neither view is the correct one; the comparison shows inconsistency, not which view
is accurate.
\item \emph{Truncation threshold}: the QC used a tail threshold of 0.30; the stricter 0.15 supported by the
truncation analysis may change the pair counts.
\item \emph{Statistics}: 43 pairs, non-normal differences for most features, and limits of agreement for the
perfusion parameters dominated by a few extreme maps; effects are reported as medians and with non-parametric
limits for that reason.
\end{itemize}

\textbf{Recommendation.}
(i) Do not pool or compare absolute AIF amplitudes, and above all absolute perfusion values, across views: for the
same injection CBV differs by about a factor two between the frontal and lateral view. Analyse each view separately,
or choose one standard view for a study and report it, together with the tube voltage.
(ii) When AIFs from different views or acquisitions must be compared, use shape-only (peak- or AUC-normalised)
features; they agree between views to within about half a frame interval.
(iii) For the AIF itself, the frontal view has less vessel overlap around the supraclinoid segment and no
saturation; the lateral view has more contrast. Together with SQ1, the frontal view with the supraclinoid AIF is the
more robust default; overlap around the mask and saturation should be checked in the lateral view.
(iv) Because the deconvolution did not cancel the view difference, the tissue side of the projection needs to be
understood before DSA perfusion values can be treated as quantitative: a matched-territory comparison of tissue
curves between views, and a phantom with known iodine at both tube voltages, are the logical next steps.
